\choiceheading

\begin{question}[index=1]
设 $\varphi(x)$ 为连续函数，$f(x)=\displaystyle\int_{-\pi}^{\pi}|\sin t-x|\varphi(t)\dd t$，现有四个命题：\par
\circlednumber{1} 若 $\varphi(x)$ 为奇函数，则 $f(x)$ 也是奇函数；\par
\circlednumber{2} 若 $\varphi(x)$ 为奇函数，则 $f(x)$ 为偶函数；\par
\circlednumber{3} 若 $\varphi(x)$ 为偶函数，则 $f(x)$ 也是偶函数；\par
\circlednumber{4} 若 $\varphi(x)$ 为偶函数，则 $f(x)$ 为奇函数。\par
以上命题中正确的是\paren
\begin{choices}[columns=4]
\item \circlednumber{1}\circlednumber{2}
\item \circlednumber{1}\circlednumber{3}
\item \circlednumber{2}\circlednumber{4}
\item \circlednumber{3}\circlednumber{4}
\end{choices}
\end{question}

\begin{question}[index=2]
若 $\displaystyle\int_{-\pi}^{\pi}(x-a_1\cos x-b_1\sin x)^2\dd x
=\min_{a,b\in\mathbb R}\left\{\int_{-\pi}^{\pi}(x-a\cos x-b\sin x)^2\dd x\right\}$，\\
则 $a_1\cos x+b_1\sin x=$\paren
\begin{choices}[columns=4]
\item $2\sin x$
\item $2\cos x$
\item $2\pi\sin x$
\item $2\pi\cos x$
\end{choices}
\end{question}

\begin{question}[index=3]
设函数 $f(x,y)$ 具有连续的偏导数，且满足 $f(x,\ln(1+x))=\dfrac{x\eu^x}{1+x}$ 及
$f(2x,x)=\eu^{2x}-\eu^x$，则
$\displaystyle\lim_{x\to0}\frac{\int_0^x f(\ln\cos t,t^2)\dd t}{x-\sin x}=$\paren
\begin{choices}[columns=4]
\item $-2$
\item $2$
\item $-1$
\item $1$
\end{choices}
\end{question}

\begin{question}[index=4]
设 $f(x,y)$ 具有二阶连续偏导数，$f'_x(x,y)>0$，$f'_y(x,y)<0$，$f''_{xy}(x,y)>0$，$D$ 为平面有界闭区域，记
\[
\begin{gathered}
I_1=\iint_D f[\eu^x,f''_{xy}(x_0,y)]\dd\sigma,\qquad
I_2=\iint_D f[x+1,f''_{yx}(x_0,y)]\dd\sigma,\\
I_3=\iint_D f[\eu^x,f'_y(x_0,y)]\dd\sigma,
\end{gathered}
\]
则\paren
\begin{choices}[columns=4]
\item $I_1>I_2>I_3$
\item $I_1>I_3>I_2$
\item $I_3>I_1>I_2$
\item $I_3>I_2>I_1$
\end{choices}
\end{question}

\begin{question}[index=5]
设 $\symbfit{A}$ 是 $m\times n$ 矩阵（$n>m$），$r(\symbfit{A})=m$，$\symbfit{B}$ 是 $n\times(n-m)$ 矩阵，
$r(\symbfit{B})=n-m$ 且 $\symbfit{A}\symbfit{B}=\symbfit{O}$。若 $\symbfit{\eta}$ 是
$\symbfit{A}\symbfit{x}=\symbf{0}$ 的解，则线性方程组 $\symbfit{B}\symbfit{y}=\symbfit{\eta}$\paren
\begin{choices}[columns=2]
\item 无解
\item 有无穷多解
\item 有唯一解
\item 解的情况不能确定
\end{choices}
\end{question}

\begin{question}[index=6]
设 $\symbfit{A}$ 为 $4$ 阶矩阵，$\symbfit{A}^*$ 为 $\symbfit{A}$ 的伴随矩阵。若
$\symbfit{A}(\symbfit{A}-\symbfit{A}^*)=\symbfit{O}$，且 $\symbfit{A}\ne\symbfit{A}^*$，\\
则 $r(\symbfit{A})$ 的取值为\paren
\begin{choices}[columns=4]
\item $0$ 或 $1$
\item $1$ 或 $3$
\item $2$ 或 $3$
\item $1$ 或 $2$
\end{choices}
\end{question}

\Needspace{7\baselineskip}
\begin{question}[index=7]
设 $\symbfit{A}$ 为 $n$ 阶正定矩阵，特征值为 $\lambda_n\ge\lambda_{n-1}\ge\cdots\ge\lambda_1>1$，
$\symbfit{Q}$ 为正交矩阵，则下列命题中，正确命题的个数是\paren\par
\circlednumber{1} $\symbfit{A}+\symbfit{Q}^{\mathrm T}\symbfit{A}^{-1}\symbfit{Q}$ 是正定矩阵。\par
\circlednumber{2} $\symbfit{A}+\symbfit{Q}^{\mathrm T}\symbfit{A}^{-1}\symbfit{Q}$ 的最小特征值为 $\lambda_1+\dfrac{1}{\lambda_1}$。\par
\circlednumber{3} $\symbfit{A}+\symbfit{Q}^{\mathrm T}\symbfit{A}^{-1}\symbfit{Q}$ 的最大特征值为 $\lambda_n+\dfrac{1}{\lambda_n}$。\par
\circlednumber{4} 若 $\symbfit{\alpha}$ 为 $\symbfit{A}$ 的特征向量，则 $\symbfit{Q}^{\mathrm T}\symbfit{\alpha}$ 为
$\symbfit{Q}^{\mathrm T}\symbfit{A}^{-1}\symbfit{Q}$ 的特征向量。
\begin{choices}[columns=4]
\item $1$
\item $2$
\item $3$
\item $4$
\end{choices}
\end{question}

\begin{question}[index=8]
设 $x_1,x_2,\ldots,x_n$ 为来自几何分布总体 $X$ 的样本观测值，$X$ 的概率分布为\\
\mbox{$P\{X=k\}=p(1-p)^{k-1}$（$k=1,2,\ldots$）}。\\
要使该样本值中等于 $2$ 的个数的期望达到最大，则 $p=$\paren
\begin{choices}[columns=4]
\item $\dfrac{1}{5}$
\item $\dfrac{1}{4}$
\item $\dfrac{1}{3}$
\item $\dfrac{1}{2}$
\end{choices}
\end{question}

\begin{question}[index=9]
设 $(X_1,X_2,\ldots,X_n)$ 和 $(Y_1,Y_2,\ldots,Y_n)$ 为来自总体 $X\sim N(\mu,\sigma^2)$ 的两个独立样本，
其样本均值分别为 $\overline X$ 和 $\overline Y$，样本方差分别为 $S_X^2$ 和 $S_Y^2$，记 $W=S_X^2+S_Y^2$。
对任意 $0<\alpha<1$，以下各式中正确的是\paren
\begin{choices}[columns=1]
\item $P\left\{\chi^2_{1-\frac{\alpha}{2}}(2n)<\dfrac{(n-1)W}{\sigma^2}<\chi^2_{\frac{\alpha}{2}}(2n)\right\}=1-\alpha$
\item $P\left\{\chi^2_{1-\frac{\alpha}{2}}(2n-2)<\dfrac{(n-1)W}{\sigma^2}<\chi^2_{\frac{\alpha}{2}}(2n-2)\right\}=1-\alpha$
\item $P\left\{-t_{\frac{\alpha}{2}}(2n)<\dfrac{\sqrt n(\overline X-\overline Y)}{W}<t_{\frac{\alpha}{2}}(2n)\right\}=1-\alpha$
\item $P\left\{-t_{\frac{\alpha}{2}}(2n-2)<\dfrac{\sqrt n(\overline X-\overline Y)}{W}<t_{\frac{\alpha}{2}}(2n-2)\right\}=\alpha$
\end{choices}
\end{question}

\begin{question}[index=10]
设总体 $X$ 的概率密度为 $f(x)=\begin{cases}
2\lambda x\eu^{-\lambda x^2}, & x>0,\\
0, & x\le 0,
\end{cases}$ $\lambda>0$ 未知。\\
对 $X$ 进行独立重复观测 $10$ 次，结果有 $4$ 次观测值大于 $1$，\\
则 $\lambda$ 的极大似然估计值 $\hat{\lambda}_L=$\paren
\begin{choices}[columns=4]
\item $\ln 5$
\item $\ln 4$
\item $\ln\dfrac{5}{2}$
\item $\ln\dfrac{5}{3}$
\end{choices}
\end{question}

\fillheading

\begin{question}[index=11]
设 $y=f(x)$ 的反函数是 $y=\varphi(x)$，且
$f(x)=\displaystyle\int_1^{2x}\eu^{t^2}\dd t+1$，则 $\varphi''(1)=$\fillin。
\end{question}

\begin{question}[index=12]
设 $y=y(x)$ 是微分方程 $y'''-y''-y'+y=0$ 的解，且满足
$\displaystyle\int_0^{+\infty}y(x)\dd x=2$，则 $y=$\fillin。
\end{question}

\begin{question}[index=13]
设数列 $\{a_n\}$ 单调减少，$\lim\limits_{n\to\infty}a_n=0$，
$S_n=\displaystyle\sum_{k=1}^n a_k$（$n=1,2,\ldots$）无界，则幂级数
$\displaystyle\sum_{n=1}^{\infty}a_n(x-1)^n$ 的收敛域为\fillin。
\end{question}

\begin{question}[index=14]
设向量场 $\symbfit{A}(x,y,z)=(z^2,x^2,y^2)$，以点 $M_0(1,1,0)$ 为圆心，$\varepsilon$ 为半径，在 $xOy$ 面上作一圆盘 $\Sigma_\varepsilon$，
面积为 $\sigma_\varepsilon$，$\Gamma$ 为该圆盘的正向边界，$\symbfit{\tau}$ 为 $\Gamma$ 上点 $(x,y,z)$ 处的单位切向量，则
$\displaystyle\lim_{\varepsilon\to0^+}\frac{1}{\sigma_\varepsilon}\oint_\Gamma\symbfit{A}\cdot\symbfit{\tau}\dd s=$\fillin。
\end{question}

\begin{question}[index=15]
设矩阵 $\symbfit{A}=\begin{bmatrix}1&2\\-2&-a\end{bmatrix}$，
$\symbfit{B}=\begin{bmatrix}1&0\\1&a\end{bmatrix}$，若
$f(x,y)=|x\symbfit{A}+y\symbfit{B}|$ 是正定二次型，\\
则 $a$ 的取值范围是\fillin。
\end{question}

\begin{question}[index=16]
设二维随机变量 $(X,Y)$ 的概率密度为 $f(x,y)=\begin{cases}
k(1-xy), & |x|\le 1,\ |y|\le 1,\\
0, & \text{其他},
\end{cases}$，\\
则 $P\!\left\{Y<\dfrac12\,\middle|\,X=\dfrac14\right\}=$\fillin。
\end{question}

\solveheading

\begin{question}[index=17,points=10,show-points=true,points-prelabel={（本题满分},points-postlabel={分）},points-separate-par=false]
设 $y=y(x)$ 是区间 $(-\pi,\pi)$ 内过点 $\left(-\dfrac{\pi}{\sqrt2},\dfrac{\pi}{\sqrt2}\right)$ 的光滑曲线。
当 $-\pi<x<0$ 时，曲线上任一点处的法线都过原点；当 $0\le x<\pi$ 时，函数 $y(x)$ 满足 $y''+y+x=0$。
求 $y(x)$ 的表达式。
\end{question}

\begin{question}[index=18,points=12,show-points=true,points-prelabel={（本题满分},points-postlabel={分）},points-separate-par=false]
设 $z=f(u,v)$ 具有二阶连续偏导数，$u=x+2y$，$v=x-3y$，且满足
$6z''_{xx}-z''_{xy}-z''_{yy}=50\lambda$，其中 $\lambda\ne 0$。若 $f'_u(u,0)=-2u$，$f(0,v)=-v^3$，求 $z=f(u,v)$ 的极值。
\end{question}

\begin{question}[index=19,points=12,show-points=true,points-prelabel={（本题满分},points-postlabel={分）},points-separate-par=false]
设 $\Omega$ 是由锥面 $x^2+(y-z)^2=(1-z)^2$（$0\le z\le 1$）与平面 $z=0$ 围成的锥体，求 $\Omega$ 的形心坐标。
\end{question}

\begin{question}[index=20,points=12,show-points=true,points-prelabel={（本题满分},points-postlabel={分）},points-separate-par=false]
设函数 $f(x)$ 在 $(-\infty,+\infty)$ 内可导，且 $f(0)=1$，$f'(x)>1$。
若 $x_1>0$，$f(x_{n+1})=f(x_n)-x_n$。
\begin{enumerate}[label=（\Roman*）,leftmargin=2.6em,labelsep=.4em,itemsep=3pt,topsep=3pt]
\item 证明 $\displaystyle\lim_{n\to\infty}x_n$ 存在，并求之；
\item 求极限 $\displaystyle\lim_{n\to\infty}(x_1+x_2+\cdots+x_n)$。
\end{enumerate}
\end{question}

\begin{question}[index=21,points=12,show-points=true,points-prelabel={（本题满分},points-postlabel={分）},points-separate-par=false]
设 $\symbfit{A}=\begin{pmatrix}1&-1&1\\a&4&b\\-3&-3&5\end{pmatrix}$，$|\symbfit{A}|=24$，
$\symbfit{\alpha}=(1,0,1)^{\mathrm T}$ 为 $\symbfit{A}^{-1}$ 的特征向量。
\begin{enumerate}[label=（\Roman*）,leftmargin=2.6em,labelsep=.4em,itemsep=3pt,topsep=3pt]
\item 求 $a,b$；
\item 求可逆矩阵 $\symbfit{P}$，使 $\symbfit{P}^{-1}\symbfit{A}\symbfit{P}$ 为对角阵；
\item 求可逆的实对称阵 $\symbfit{Q}$，使 $\symbfit{Q}^{-1}\symbfit{A}\symbfit{Q}=\symbfit{A}^{\mathrm T}$。
\end{enumerate}
\end{question}

\begin{question}[index=22,points=12,show-points=true,points-prelabel={（本题满分},points-postlabel={分）},points-separate-par=false]
设随机变量 $X,Y$ 相互独立，$X\sim P(\lambda)$，$Y\sim P(\mu)$，$\lambda>0$，$\mu>0$，令 $Z=X+Y$。
\begin{enumerate}[label=（\arabic*）,leftmargin=2.6em,labelsep=.4em,itemsep=3pt,topsep=3pt]
\item 证明：在条件 $Z=n$（$n=0,1,2,\ldots$）下，$X$ 服从二项分布，并写出该二项分布的\mbox{两个参数}；
\item 求协方差 $\operatorname{Cov}(X,Z)$ 与相关系数 $\rho_{XZ}$；
\item 取 $\lambda=2$，$\mu=3$，计算条件概率 $P(X=1\mid Z=4)$。
\end{enumerate}
\end{question}

